\documentclass[12pt,a4paper]{exam}
\usepackage[utf8]{inputenc}
\usepackage{amsmath,amssymb,amsfonts}
\usepackage{geometry}
\usepackage{enumitem}
\usepackage{graphicx}
\usepackage{hyperref}
\usepackage{xcolor}
\geometry{a4paper, top=2cm, bottom=2cm, left=2.5cm, right=2.5cm}
\hypersetup{colorlinks=true, linkcolor=blue, urlcolor=blue}
\renewcommand{\thequestion}{\textbf{\arabic{question}}}
\renewcommand{\questionlabel}{\thequestion.}
\pointname{ marks}
\pointsdroppedatright
\bracketedpoints
\setlength{\parindent}{0pt}
\setlength{\emergencystretch}{3em}
\allowdisplaybreaks
\newsavebox{\schmathbox}
\newcommand{\fitmath}[1]{%
  \sbox{\schmathbox}{$\displaystyle #1$}%
  \ifdim\wd\schmathbox>\linewidth
    \resizebox{\linewidth}{!}{\usebox{\schmathbox}}%
  \else
    \usebox{\schmathbox}%
  \fi}

\begin{document}

\thispagestyle{empty}
\begin{center}
  \textbf{\LARGE Diag Solutions} \\[0.3cm]
  \rule{\textwidth}{0.5pt}
\end{center}

\vspace{0.3cm}

\begin{questions}
\question \textbf{1.} In $\triangle PQR$, $S$ and $T$ are points on sides $PQ$ and $PR$ respectively such that $\angle PST = \angle PRQ$. If $PS = 2\text{ cm}$, $SQ = 3\text{ cm}$, and $PT = 2.5\text{ cm}$, find $TR$.

\textbf{Answer:}
\begin{center}
\fitmath{1.5\text{ cm}}
\end{center}

\textbf{Solution:}
\begin{enumerate}
  \item Show $\triangle PST \sim \triangle PRQ$ by AA similarity criterion.
  \item Write the ratio of sides and solve for $TR$.
\end{enumerate}
\textbf{Derivation:}
\begin{center}
\fitmath{\begin{aligned}
\triangle PST \sim \triangle PRQ \\
\implies \frac{PS}{PR} = \frac{PT}{PQ} \text{ is incorrect; use } \frac{PS}{PR} = \frac{PT}{PQ} \text{ is not needed, use } \frac{PS}{PT} = \frac{PT}{PR} \text{ is not right; use } \frac{PS}{PR} = \frac{PT}{PQ} \\
\implies \frac{PS}{PT+TR} = \frac{PT}{PS+SQ} \\
\implies \frac{2}{2.5+TR} = \frac{2.5}{5} \\
\implies 10 = 6.25 + 2.5TR \\
\implies 3.75 = 2.5TR \\
\implies TR = 1.5\text{ cm}
\end{aligned}}
\end{center}
\hfill \textit{Triangles — Similarity of Triangles}
\vspace{0.3cm}

\question \textbf{2.} Diagonals of a trapezium $\text{ABCD}$ with $AB \parallel DC$ intersect each other at the point $O$. If $AB = 2CD$, find the ratio of the areas of triangles $\triangle AOB$ and $\triangle COD$.

\textbf{Answer:}
\begin{center}
\fitmath{4:1}
\end{center}

\textbf{Solution:}
\begin{enumerate}
  \item Prove $\triangle AOB \sim \triangle COD$ by AA similarity.
  \item Use the ratio of areas theorem.
\end{enumerate}
\textbf{Derivation:}
\begin{center}
\fitmath{\begin{aligned}
\triangle AOB \sim \triangle COD \\
\implies \frac{\text{Area}(AOB)}{\text{Area}(COD)} = \left(\frac{AB}{CD}\right)^2 = \left(\frac{2CD}{CD}\right)^2 = 2^2 = 4
\end{aligned}}
\end{center}
\hfill \textit{Triangles — Areas of Similar Triangles}
\vspace{0.3cm}

\question \textbf{3.} In $\triangle ABC$, $DE \parallel BC$ such that $AD = (4x - 3)$ cm, $AE = (8x - 7)$ cm, $BD = (3x - 1)$ cm, and $EC = (5x - 3)$ cm. Find the value of $x$.

\textbf{Answer:}
\begin{center}
\fitmath{x = 1}
\end{center}

\textbf{Solution:}
\begin{enumerate}
  \item Since $DE \parallel BC$, by Basic Proportionality Theorem (BPT), $\frac{AD}{DB} = \frac{AE}{EC}$.
  \item Substitute the given values into the equation: $\frac{4x-3}{3x-1} = \frac{8x-7}{5x-3}$.
  \item Cross-multiply to get: $(4x-3)(5x-3) = (8x-7)(3x-1)$.
  \item Solve the resulting linear equation to find $x$.
\end{enumerate}
\textbf{Derivation:}
\begin{center}
\fitmath{\begin{aligned}
20x^2 - 12x - 15x + 9 = 24x^2 - 8x - 21x + 7 \\
\implies 20x^2 - 27x + 9 = 24x^2 - 29x + 7 \\
\implies 4x^2 - 2x - 2 = 0 \\
\implies 2x^2 - x - 1 = 0 \\
\implies (2x+1)(x-1) = 0. \text{ Since } x > 1, x = 1.
\end{aligned}}
\end{center}
\hfill \textit{Triangles — Basic Proportionality Theorem}
\vspace{0.3cm}

\question \textbf{4.} Two poles of heights $6$ m and $11$ m stand on a plane ground. If the distance between their feet is $12$ m, find the distance between their tops.

\textbf{Answer:}
\begin{center}
\fitmath{13 m}
\end{center}

\textbf{Solution:}
\begin{enumerate}
  \item Let $AB = 6$ m and $CD = 11$ m be the two poles. The distance $BD = 12$ m.
  \item Draw a line from $A$ parallel to $BD$ meeting $CD$ at $E$. Then $CE = CD - ED = 11 - 6 = 5$ m.
  \item In right-angled triangle $\triangle AEC$, use the Pythagorean theorem to find the hypotenuse $AC$.
  \item $AC^2 = AE^2 + CE^2$.
\end{enumerate}
\textbf{Derivation:}
\begin{center}
\fitmath{AE = BD = 12 \text{ m}. CE = 11 - 6 = 5 \text{ m}. AC = \sqrt{12^2 + 5^2} = \sqrt{144 + 25} = \sqrt{169} = 13 \text{ m}.}
\end{center}
\hfill \textit{Triangles — Pythagoras Theorem}
\vspace{0.3cm}

\question \textbf{5.} In $\triangle PQR$, $S$ and $T$ are points on sides $PQ$ and $PR$ respectively such that $\angle PST = \angle PRQ$. Show that $\triangle PQR \sim \triangle PTS$.

\textbf{Answer:}
\begin{center}
\fitmath{\text{Proof}}
\end{center}

\textbf{Solution:}
\begin{enumerate}
  \item Identify the common angle in both triangles: $\angle QPR = \angle TPS$ (Common).
  \item Use the given condition $\angle PST = \angle PRQ$.
  \item Conclude similarity by AA (Angle-Angle) criterion.
\end{enumerate}
\textbf{Derivation:}
\begin{center}
\fitmath{\text{In } \triangle PQR \text{ and } \triangle PTS: \angle P = \angle P \text{ (Common)}, \angle PST = \angle PRQ \text{ (Given)}. \text{ Therefore, } \triangle PQR \sim \triangle PTS \text{ by AA similarity criterion.}}
\end{center}
\hfill \textit{Triangles — Criteria for Similarity of Triangles}
\vspace{0.3cm}

\question \textbf{6.} The areas of two similar triangles are $121$ cm$^2$ and $64$ cm$^2$ respectively. If the median of the first triangle is $12.1$ cm, find the corresponding median of the other triangle.

\textbf{Answer:}
\begin{center}
\fitmath{8.8 cm}
\end{center}

\textbf{Solution:}
\begin{enumerate}
  \item Use the theorem: Ratio of areas of similar triangles is equal to the ratio of the squares of their corresponding medians.
  \item Set up the equation: $\frac{\text{Area}_1}{\text{Area}_2} = \left(\frac{m_1}{m_2}\right)^2$.
  \item Substitute values: $\frac{121}{64} = \left(\frac{12.1}{m_2}\right)^2$.
  \item Take the square root and solve for $m_2$.
\end{enumerate}
\textbf{Derivation:}
\begin{center}
\fitmath{\begin{aligned}
\frac{11}{8} = \frac{12.1}{m_2} \\
\implies m_2 = \frac{12.1 \times 8}{11} = 1.1 \times 8 = 8.8 \text{ cm}.
\end{aligned}}
\end{center}
\hfill \textit{Triangles — Areas of Similar Triangles}
\vspace{0.3cm}

\question \textbf{7.} In $\triangle ABC$, $AD \perp BC$ such that $AD^2 = BD \cdot CD$. Prove that $\angle BAC = 90^\circ$.

\textbf{Answer:}
\begin{center}
\fitmath{\text{Proof}}
\end{center}

\textbf{Solution:}
\begin{enumerate}
  \item Since $AD \perp BC$, we have two right triangles $\triangle ADB$ and $\triangle ADC$.
  \item Express the sides using Pythagoras: $AB^2 = AD^2 + BD^2$ and $AC^2 = AD^2 + CD^2$.
  \item Add them and use the given condition $AD^2 = BD \cdot CD$ to show $AB^2 + AC^2 = BC^2$.
\end{enumerate}
\textbf{Derivation:}
\begin{center}
\fitmath{AB^2 + AC^2 = AD^2 + BD^2 + AD^2 + CD^2 = 2AD^2 + BD^2 + CD^2. \text{ Substitute } AD^2=BD \cdot CD: AB^2 + AC^2 = 2(BD \cdot CD) + BD^2 + CD^2 = (BD + CD)^2 = BC^2. \text{ By converse of Pythagoras, } \angle BAC = 90^\circ.}
\end{center}
\hfill \textit{Triangles — Pythagoras Theorem}
\vspace{0.3cm}

\question \textbf{8.} State and prove the Basic Proportionality Theorem (Thales Theorem).

\textbf{Answer:}
\begin{center}
\fitmath{\text{If a line is drawn parallel to one side of a triangle to intersect the other two sides in distinct points}, \text{the other two sides are divided in the same ratio}.}
\end{center}

\textbf{Solution:}
\begin{enumerate}
  \item Draw a triangle ABC with a line DE parallel to BC intersecting AB at D and AC at E.
  \item Construct perpendiculars DM on AC and EN on AB. Join BE and CD.
  \item Calculate the area of triangles ADE, BDE, and CDE.
  \item Use the property that triangles on the same base and between the same parallels have equal areas.
  \item Divide the ratio of areas to prove the theorem.
\end{enumerate}
\textbf{Derivation:}
In  $\triangle$ ADE  and  $\triangle$ BDE ,  $\frac{ar(ADE)}{ar(BDE)} = \frac{\frac{1}{2} \cdot AD \cdot EN}{\frac{1}{2} \cdot BD \cdot EN} = \frac{AD}{BD}$ . Similarly,  $\frac{ar(ADE)}{ar(CDE)} = \frac{AE}{CE}$ . Since  ar(BDE) = ar(CDE)  (same base  DE  and same parallels  DE $\parallel$ BC ), it follows that  $\frac{AD}{BD} = \frac{AE}{CE}$ .
\hfill \textit{Triangles — Basic Proportionality Theorem}
\vspace{0.3cm}

\question \textbf{9.} Prove that the ratio of the areas of two similar triangles is equal to the square of the ratio of their corresponding sides.

\textbf{Answer:}
\begin{center}
\fitmath{\frac{ar(\triangle ABC)}{ar(\triangle PQR)} = (\frac{AB}{PQ})^2 = (\frac{BC}{QR})^2 = (\frac{AC}{PR})^2}
\end{center}

\textbf{Solution:}
\begin{enumerate}
  \item Draw two similar triangles ABC and PQR.
  \item Draw altitudes AM $\perp$ BC and PN $\perp$ QR.
  \item Express the ratio of areas as $\frac{\frac{1}{2} \cdot BC \cdot AM}{\frac{1}{2} \cdot QR \cdot PN}$.
  \item Prove $\triangle ABM \sim \triangle PQN$ to show $\frac{AM}{PN} = \frac{AB}{PQ}$.
  \item Substitute this ratio back into the area equation to get the square of the ratio.
\end{enumerate}
\textbf{Derivation:}
 $\frac{ar(ABC)}{ar(PQR)} = \frac{BC \cdot AM}{QR \cdot PN} = \frac{BC}{QR} \cdot \frac{AM}{PN}$ . Since  $\triangle$ ABM $\sim \triangle$ PQN ,  $\frac{AM}{PN} = \frac{AB}{PQ}$ . Also  $\frac{AB}{PQ} = \frac{BC}{QR}  ($given similarity). Thus,  $\frac{ar(ABC)}{ar(PQR)} = \frac{BC}{QR} \cdot \frac{BC}{QR} = (\frac{BC}{QR})^2$ .
\hfill \textit{Triangles — Areas of Similar Triangles}
\vspace{0.3cm}

\question \textbf{10.} In $\triangle ABC$, $AD \perp BC$ such that $AD^2 = BD \cdot CD$. Prove that $\triangle ABC$ is a right-angled triangle at $A$.

\textbf{Answer:}
\begin{center}
\fitmath{\angle BAC = 90^\circ}
\end{center}

\textbf{Solution:}
\begin{enumerate}
  \item Observe the triangle split into two smaller triangles $\triangle ABD$ and $\triangle ACD$.
  \item Write the ratio $\frac{AD}{BD} = \frac{CD}{AD}$.
  \item Include the angle $\angle ADB = \angle ADC = 90^\circ$.
  \item Prove $\triangle ABD \sim \triangle CAD$ by SAS similarity.
  \item Equate corresponding angles $\angle BAD = \angle ACD$ and $\angle CAD = \angle ABD$.
  \item Use the angle sum property of $\triangle ABC$ to show $\angle BAC = 90^\circ$.
\end{enumerate}
\textbf{Derivation:}
Given  AD\textasciicircum{}2 = BD $\cdot$ CD $\implies \frac{AD}{BD} = \frac{CD}{AD}$ . In  $\triangle$ ABD  and  $\triangle$ CAD ,  $\angle$ ADB = $\angle$ ADC = 90\textasciicircum{}$\circ$ . By SAS similarity,  $\triangle$ ABD $\sim \triangle$ CAD . Thus  $\angle$ BAD = $\angle C$  and  $\angle$ CAD = $\angle B$ . In  $\triangle$ ABC ,  $\angle A + \angle B + \angle C = 180^\circ$ . Substituting gives  $\angle A + \angle$ CAD + $\angle$ BAD = 180\textasciicircum{}$\circ \implies 2\angle A = 180^\circ \implies \angle A = 90^\circ$ .
\hfill \textit{Triangles — Similarity of Triangles}
\vspace{0.3cm}

\question \textbf{11.} A vertical stick $2$ m long casts a shadow $1$ m long on the ground. At the same time, a tower casts a shadow $5$ m long. Find the height of the tower using the properties of similar triangles.

\textbf{Answer:}
\begin{center}
\fitmath{10 \text{meters}}
\end{center}

\textbf{Solution:}
\begin{enumerate}
  \item Identify two triangles: one formed by the stick and its shadow, the other by the tower and its shadow.
  \item State that at the same time, the angle of elevation of the sun is the same for both objects.
  \item Identify that both triangles are right-angled and share an equal angle, thus they are similar by AA criterion.
  \item Set up the proportion: $\frac{\text{Height of Stick}}{\text{Shadow of Stick}} = \frac{\text{Height of Tower}}{\text{Shadow of Tower}}$.
  \item Substitute the given values and solve for the height of the tower.
\end{enumerate}
\textbf{Derivation:}
Let height of tower be  h .  $\triangle$ ABC $\sim \triangle$ PQR  (where  ABC  is stick triangle and  PQR  is tower triangle). Therefore,  $\frac{2}{1} = \frac{h}{5}$ .  h = 2 $\times 5 = 10  m.$
\hfill \textit{Triangles — Applications of Similar Triangles}
\vspace{0.3cm}

\question \textbf{12.} The distance of the point $P(3, 4)$ from the origin is:

\textbf{Answer:}
(C)

\textbf{Solution:}
\begin{enumerate}
  \item Use the distance formula from the origin $d = \sqrt{x^2 + y^2}$
  \item Substitute $x=3$ and $y=4$
  \item Calculate $\sqrt{3^2 + 4^2} = \sqrt{9 + 16} = \sqrt{25} = 5$
\end{enumerate}
\textbf{Derivation:}
\begin{center}
\fitmath{d = \sqrt{3^2 + 4^2} = \sqrt{9+16} = \sqrt{25} = 5}
\end{center}
\hfill \textit{Coordinate Geometry — Distance Formula}
\vspace{0.3cm}

\question \textbf{13.} The coordinates of the midpoint of the line segment joining the points $A(2, 4)$ and $B(4, 6)$ are:

\textbf{Answer:}
(B)

\textbf{Solution:}
\begin{enumerate}
  \item Use the midpoint formula $M = (\frac{x_1+x_2}{2}, \frac{y_1+y_2}{2})$
  \item Substitute the values: $x_1=2, x_2=4, y_1=4, y_2=6$
  \item Calculate: $x = \frac{2+4}{2} = 3$ and $y = \frac{4+6}{2} = 5$
\end{enumerate}
\textbf{Derivation:}
\begin{center}
\fitmath{M = (\frac{2+4}{2}, \frac{4+6}{2}) = (3, 5)}
\end{center}
\hfill \textit{Coordinate Geometry — Section Formula}
\vspace{0.3cm}

\question \textbf{14.} The abscissa of the point $(5, 8)$ is:

\textbf{Answer:}
(A)

\textbf{Solution:}
\begin{enumerate}
  \item Recall that in a coordinate point $(x, y)$, $x$ is the abscissa and $y$ is the ordinate.
  \item For the point $(5, 8)$, the $x$-coordinate is 5.
\end{enumerate}
\textbf{Derivation:}
\begin{center}
\fitmath{\text{Abscissa} = x = 5}
\end{center}
\hfill \textit{Coordinate Geometry — Introduction to Cartesian Plane}
\vspace{0.3cm}

\question \textbf{15.} If the distance between points $(2, -2)$ and $(-1, x)$ is 5, then the value of $x$ is not asked, but what is the distance formula expression?

\textbf{Answer:}
(D)

\textbf{Solution:}
\begin{enumerate}
  \item The distance formula between two points $(x_1, y_1)$ and $(x_2, y_2)$ is $\sqrt{(x_2-x_1)^2 + (y_2-y_1)^2}$
  \item Substitute $x_1=2, y_1=-2, x_2=-1, y_2=x$
\end{enumerate}
\textbf{Derivation:}
\begin{center}
\fitmath{d = \sqrt{(-1-2)^2 + (x - (-2))^2} = \sqrt{(-3)^2 + (x+2)^2}}
\end{center}
\hfill \textit{Coordinate Geometry — Distance Formula}
\vspace{0.3cm}

\question \textbf{16.} The point which lies on the y-axis is:

\textbf{Answer:}
(C)

\textbf{Solution:}
\begin{enumerate}
  \item A point on the y-axis has its $x$-coordinate equal to 0.
  \item Check the options for $x=0$.
\end{enumerate}
\textbf{Derivation:}
Point on y-axis is of the form (0, y).
\hfill \textit{Coordinate Geometry — Introduction to Cartesian Plane}
\vspace{0.3cm}

\question \textbf{17.} Assertion (A): The point $(0, 5)$ lies on the y-axis. Reason (R): For any point on the y-axis, the x-coordinate is zero.

\textbf{Answer:}
(A)

\textbf{Solution:}
\begin{enumerate}
  \item Check the definition of y-axis points.
  \item Verify if (0, 5) satisfies the condition x=0.
  \item Conclude that R explains A.
\end{enumerate}
\textbf{Derivation:}
A point  (x, y)  lies on the y-axis if and only if  x=0 . Since  (0, 5)  has  x=0 , it lies on the y-axis.
\hfill \textit{Coordinate Geometry — Introduction to Coordinate Geometry}
\vspace{0.3cm}

\question \textbf{18.} Assertion (A): The distance of the point $P(3, 4)$ from the origin $(0, 0)$ is $5$ units. Reason (R): The distance of a point $(x, y)$ from the origin is calculated by $\sqrt{x^2 + y^2}$.

\textbf{Answer:}
(A)

\textbf{Solution:}
\begin{enumerate}
  \item Apply the distance formula from the origin.
  \item Substitute x=3 and y=4.
  \item Calculate $\sqrt{3^2 + 4^2} = \sqrt{9+16} = \sqrt{25} = 5$.
\end{enumerate}
\textbf{Derivation:}
\begin{center}
\fitmath{d = \sqrt{(3-0)^2 + (4-0)^2} = \sqrt{3^2 + 4^2} = \sqrt{9+16} = 5}
\end{center}
\hfill \textit{Coordinate Geometry — Distance Formula}
\vspace{0.3cm}

\question \textbf{19.} Assertion (A): The midpoint of the line segment joining $(2, 4)$ and $(6, 8)$ is $(4, 6)$. Reason (R): The midpoint formula is $(\frac{x_1+x_2}{2}, \frac{y_1+y_2}{2})$.

\textbf{Answer:}
(A)

\textbf{Solution:}
\begin{enumerate}
  \item Identify coordinates $(x_1, y_1) = (2, 4)$ and $(x_2, y_2) = (6, 8)$.
  \item Use the midpoint formula.
  \item Calculate $(\frac{2+6}{2}, \frac{4+8}{2}) = (4, 6)$.
\end{enumerate}
\textbf{Derivation:}
\begin{center}
\fitmath{M = (\frac{2+6}{2}, \frac{4+8}{2}) = (\frac{8}{2}, \frac{12}{2}) = (4, 6)}
\end{center}
\hfill \textit{Coordinate Geometry — Section Formula/Midpoint}
\vspace{0.3cm}

\question \textbf{20.} Assertion (A): The points $(1, 1), (2, 2)$ and $(3, 3)$ are collinear. Reason (R): Three points are collinear if the sum of distances between any two pairs is equal to the distance between the third pair.

\textbf{Answer:}
(B)

\textbf{Solution:}
\begin{enumerate}
  \item Check the slopes: slope between (1,1) and (2,2) is (2-1)/(2-1) = 1.
  \item Slope between (2,2) and (3,3) is (3-2)/(3-2) = 1.
  \item Since slopes are equal, they are collinear.
  \item R is a true statement, but typically collinearity is defined by slope or area = 0 in Class 10.
\end{enumerate}
\textbf{Derivation:}
Slope\_\{AB\} = 1, Slope\_\{BC\} = 1. Since slopes are equal, points lie on the same line.
\hfill \textit{Coordinate Geometry — Collinearity}
\vspace{0.3cm}

\question \textbf{21.} Assertion (A): The area of a triangle with vertices $(0, 0), (4, 0),$ and $(0, 3)$ is $6$ sq units. Reason (R): Area of a triangle with coordinates $(0, 0), (x_1, y_1), (x_2, y_2)$ is $\frac{1}{2}|x_1y_2 - x_2y_1|$.

\textbf{Answer:}
(A)

\textbf{Solution:}
\begin{enumerate}
  \item Use the simplified area formula for a triangle with one vertex at origin.
  \item Substitute $(x_1, y_1) = (4, 0)$ and $(x_2, y_2) = (0, 3)$.
  \item Calculate $\frac{1}{2}|(4)(3) - (0)(0)| = \frac{1}{2}(12) = 6$.
\end{enumerate}
\textbf{Derivation:}
\begin{center}
\fitmath{\text{Area} = \frac{1}{2}|4 \times 3 - 0 \times 0| = 6}
\end{center}
\hfill \textit{Coordinate Geometry — Area of Triangle}
\vspace{0.3cm}

\question \textbf{22.} Find the ratio in which the line segment joining the points $A(-3, 10)$ and $B(6, -8)$ is divided by the point $P(-1, 6)$.

\textbf{Answer:}
\begin{center}
\fitmath{2:7}
\end{center}

\textbf{Solution:}
\begin{enumerate}
  \item Use the section formula: $P(x, y) = \left( \frac{mx_2 + nx_1}{m+n}, \frac{my_2 + ny_1}{m+n} \right)$
  \item Substitute the coordinates and assume ratio $k:1$
  \item Solve for $k$ using the x-coordinate equation
  \item Verify with the y-coordinate
\end{enumerate}
\textbf{Derivation:}
\begin{center}
\fitmath{\begin{aligned}
-1 = \frac{k(6) + 1(-3)}{k+1} \\
\implies -k - 1 = 6k - 3 \\
\implies 2 = 7k \\
\implies k = \frac{2}{7}
\end{aligned}}
\end{center}
\hfill \textit{Coordinate Geometry — Section Formula}
\vspace{0.3cm}

\question \textbf{23.} If the points $A(6, 1)$, $B(8, 2)$, $C(9, 4)$, and $D(p, 3)$ are the vertices of a parallelogram taken in order, find the value of $p$.

\textbf{Answer:}
\begin{center}
\fitmath{p = 7}
\end{center}

\textbf{Solution:}
\begin{enumerate}
  \item Recall that diagonals of a parallelogram bisect each other.
  \item Find the midpoint of diagonal AC.
  \item Find the midpoint of diagonal BD.
  \item Equate the midpoints and solve for $p$.
\end{enumerate}
\textbf{Derivation:}
Midpoint of AC = ($\frac{6+9}{2}, \frac{1+4}{2}) = (7.5, 2.5)$; Midpoint of BD = ($\frac{8+p}{2}, \frac{2+3}{2}) = (\frac{8+p}{2}, 2.5)$; $\frac{8+p}{2} = 7.5 \implies 8+p = 15 \implies p = 7$
\hfill \textit{Coordinate Geometry — Midpoint Formula}
\vspace{0.3cm}

\question \textbf{24.} Find the distance between the points $P(a+b, b-a)$ and $Q(a-b, a+b)$.

\textbf{Answer:}
\begin{center}
\fitmath{2\sqrt{a^2+b^2}}
\end{center}

\textbf{Solution:}
\begin{enumerate}
  \item Apply the distance formula: $d = \sqrt{(x_2-x_1)^2 + (y_2-y_1)^2}$
  \item Substitute the given coordinates.
  \item Simplify the algebraic terms inside the square root.
  \item Calculate the square root.
\end{enumerate}
\textbf{Derivation:}
\begin{center}
\fitmath{d = \sqrt{((a-b)-(a+b))^2 + ((a+b)-(b-a))^2} = \sqrt{(-2b)^2 + (2a)^2} = \sqrt{4b^2 + 4a^2} = 2\sqrt{a^2+b^2}}
\end{center}
\hfill \textit{Coordinate Geometry — Distance Formula}
\vspace{0.3cm}

\question \textbf{25.} Find the value of $y$ for which the distance between the points $P(2, -3)$ and $Q(10, y)$ is 10 units.

\textbf{Answer:}
\begin{center}
\fitmath{y = 3 \text{or y} = -9}
\end{center}

\textbf{Solution:}
\begin{enumerate}
  \item Use the distance formula $PQ^2 = (x_2-x_1)^2 + (y_2-y_1)^2$.
  \item Substitute the known values into the equation.
  \item Solve the resulting quadratic equation for $y$.
\end{enumerate}
\textbf{Derivation:}
\begin{center}
\fitmath{\begin{aligned}
10^2 = (10-2)^2 + (y-(-3))^2 \\
\implies 100 = 64 + (y+3)^2 \\
\implies 36 = (y+3)^2 \\
\implies y+3 = \pm 6 \\
\implies y = 3, -9
\end{aligned}}
\end{center}
\hfill \textit{Coordinate Geometry — Distance Formula}
\vspace{0.3cm}

\question \textbf{26.} Find the coordinates of a point $A$, where $AB$ is the diameter of a circle whose centre is $(2, -3)$ and $B$ is $(1, 4)$.

\textbf{Answer:}
\begin{center}
\fitmath{A(3, -10)}
\end{center}

\textbf{Solution:}
\begin{enumerate}
  \item Let $A$ be $(x, y)$. The centre $(2, -3)$ is the midpoint of diameter $AB$.
  \item Use the midpoint formula: $\frac{x+1}{2} = 2$ and $\frac{y+4}{2} = -3$.
  \item Solve for $x$ and $y$.
\end{enumerate}
\textbf{Derivation:}
\begin{center}
\fitmath{\begin{aligned}
x+1 = 4 \\
\implies x = 3; y+4 = -6 \\
\implies y = -10. \text{Point A is } (3, -10)
\end{aligned}}
\end{center}
\hfill \textit{Coordinate Geometry — Midpoint Formula}
\vspace{0.3cm}

\question \textbf{27.} Find the ratio in which the point $P(x, 2)$ divides the line segment joining the points $A(12, 5)$ and $B(4, -3)$. Also, find the value of $x$.

\textbf{Answer:}
\begin{center}
\fitmath{\text{Ratio is} 3:5, x = 9}
\end{center}

\textbf{Solution:}
\begin{enumerate}
  \item Use the section formula for the y-coordinate to find the ratio k:1.
  \item Set the y-coordinate formula equal to 2.
  \item Solve for k.
  \item Use the value of k to find the x-coordinate using the section formula.
  \item State the final ratio and the value of x.
\end{enumerate}
\textbf{Derivation:}
Let the ratio be  k:1 . The y-coordinate is given by  y = $\frac{m_1y_2 + m_2y_1}{m_1 + m_2}$ . Substituting values:  2 = $\frac{k(-3) + 1(5)}{k + 1} \implies 2k + 2 = -3k + 5 \implies 5k = 3 \implies k = 3/5$ . Thus, the ratio is  3:5 . For x:  x = $\frac{3(4) + 5(12)}{3 + 5} = \frac{12 + 60}{8} = \frac{72}{8} = 9$ .
\hfill \textit{Coordinate Geometry — Section Formula}
\vspace{0.3cm}

\question \textbf{28.} If the points $A(6, 1)$, $B(8, 2)$, $C(9, 4)$, and $D(p, 3)$ are the vertices of a parallelogram taken in order, find the value of $p$.

\textbf{Answer:}
\begin{center}
\fitmath{p = 7}
\end{center}

\textbf{Solution:}
\begin{enumerate}
  \item Recall the property that diagonals of a parallelogram bisect each other.
  \item Find the midpoint of diagonal AC.
  \item Find the midpoint of diagonal BD in terms of p.
  \item Equate the midpoints and solve for p.
\end{enumerate}
\textbf{Derivation:}
Midpoint of AC =  ($\frac{6+9}{2}, \frac{1+4}{2}) = (7.5, 2.5)$ . Midpoint of BD =  ($\frac{8+p}{2}, \frac{2+3}{2}) = (\frac{8+p}{2}, 2.5)$ . Equating x-coordinates:  7.5 = $\frac{8+p}{2} \implies 15 = 8 + p \implies p = 7$ .
\hfill \textit{Coordinate Geometry — Midpoint Formula}
\vspace{0.3cm}

\question \textbf{29.} Find the area of a quadrilateral whose vertices, taken in order, are $(-4, -2)$, $(-3, -5)$, $(3, -2)$, and $(2, 3)$.

\textbf{Answer:}
\begin{center}
\fitmath{28 sq. \text{units}}
\end{center}

\textbf{Solution:}
\begin{enumerate}
  \item Divide the quadrilateral into two triangles by drawing a diagonal.
  \item Use the area of triangle formula: $\frac{1}{2}|x_1(y_2 - y_3) + x_2(y_3 - y_1) + x_3(y_1 - y_2)|$.
  \item Calculate the area of triangle 1.
  \item Calculate the area of triangle 2.
  \item Add the two areas to get the total area.
\end{enumerate}
\textbf{Derivation:}
Area of  $\triangle$ ABC  with vertices  (-4, -2), (-3, -5), (3, -2) :  $\frac{1}{2}|-4(-5 - (-2)) + (-3)(-2 - (-2)) + 3(-2 - (-5))| = \frac{1}{2}|12 + 0 + 9| = 10.5$ . Area of  $\triangle$ ADC  with vertices  (-4, -2), (2, 3), (3, -2) :  $\frac{1}{2}|-4(3 - (-2)) + 2(-2 - (-2)) + 3(-2 - 3)| = \frac{1}{2}|-20 + 0 - 15| = 17.5$ . Total Area =  10.5 + 17.5 = 28 .
\hfill \textit{Coordinate Geometry — Area of Triangle}
\vspace{0.3cm}

\question \textbf{30.} Find the coordinates of the points which divide the line segment joining $A(-2, 2)$ and $B(2, 8)$ into four equal parts.

\textbf{Answer:}
\begin{center}
\fitmath{(-1, 3.5), (0, 5), (1, 6.5)}
\end{center}

\textbf{Solution:}
\begin{enumerate}
  \item Identify the midpoints needed to divide the segment into 4 parts.
  \item Find the midpoint M of AB.
  \item Find the midpoint of AM to get the first point.
  \item Find the midpoint of MB to get the third point.
  \item List all three points.
\end{enumerate}
\textbf{Derivation:}
Midpoint  M  of  AB = ($\frac{-2+2}{2}, \frac{2+8}{2}) = (0, 5)$ . Midpoint  P  of  AM = ($\frac{-2+0}{2}, \frac{2+5}{2}) = (-1, 3.5)$ . Midpoint  Q  of  MB = ($\frac{0+2}{2}, \frac{5+8}{2}) = (1, 6.5)$ . The points are  (-1, 3.5), (0, 5), (1, 6.5) .
\hfill \textit{Coordinate Geometry — Section Formula Application}
\vspace{0.3cm}

\question \textbf{31.} If $\sin \theta = \frac{3}{5}$, then the value of $\cos \theta$ is:

\textbf{Answer:}
(C)

\textbf{Solution:}
\begin{enumerate}
  \item Use the identity $\sin^2 \theta + \cos^2 \theta = 1$
  \item Substitute the value of $\sin \theta$
  \item Solve for $\cos \theta$
\end{enumerate}
\textbf{Derivation:}
\begin{center}
\fitmath{\begin{aligned}
\cos^2 \theta = 1 - \sin^2 \theta = 1 - (\frac{3}{5})^2 = 1 - \frac{9}{25} = \frac{16}{25} \\
\implies \cos \theta = \frac{4}{5}
\end{aligned}}
\end{center}
\hfill \textit{Introduction to Trigonometry — Trigonometric Identities}
\vspace{0.3cm}

\question \textbf{32.} The value of $\tan 45^\circ + \cot 45^\circ$ is:

\textbf{Answer:}
(B)

\textbf{Solution:}
\begin{enumerate}
  \item Recall the trigonometric values for standard angles.
  \item Substitute $\tan 45^\circ = 1$ and $\cot 45^\circ = 1$.
  \item Add the values.
\end{enumerate}
\textbf{Derivation:}
\begin{center}
\fitmath{\tan 45^\circ + \cot 45^\circ = 1 + 1 = 2}
\end{center}
\hfill \textit{Introduction to Trigonometry — Trigonometric Ratios of Specific Angles}
\vspace{0.3cm}

\question \textbf{33.} If $3\tan \theta = 3$, then $\theta$ is equal to:

\textbf{Answer:}
(C)

\textbf{Solution:}
\begin{enumerate}
  \item Simplify the equation to find $\tan \theta$.
  \item Identify the angle $\theta$ for which $\tan \theta = 1$.
\end{enumerate}
\textbf{Derivation:}
\begin{center}
\fitmath{\begin{aligned}
3\tan \theta = 3 \\
\implies \tan \theta = 1 \\
\implies \theta = 45^\circ
\end{aligned}}
\end{center}
\hfill \textit{Introduction to Trigonometry — Trigonometric Ratios}
\vspace{0.3cm}

\question \textbf{34.} The value of $\frac{\sin 18^\circ}{\cos 72^\circ}$ is:

\textbf{Answer:}
(A)

\textbf{Solution:}
\begin{enumerate}
  \item Use the complementary angle identity $\cos(90^\circ - \theta) = \sin \theta$.
  \item Substitute $\cos 72^\circ = \cos(90^\circ - 18^\circ) = \sin 18^\circ$.
  \item Divide the terms.
\end{enumerate}
\textbf{Derivation:}
\begin{center}
\fitmath{\frac{\sin 18^\circ}{\cos 72^\circ} = \frac{\sin 18^\circ}{\sin 18^\circ} = 1}
\end{center}
\hfill \textit{Introduction to Trigonometry — Trigonometric Ratios of Complementary Angles}
\vspace{0.3cm}

\question \textbf{35.} Which of the following is equal to $\sec^2 \theta - 1$?

\textbf{Answer:}
(D)

\textbf{Solution:}
\begin{enumerate}
  \item Recall the fundamental trigonometric identity $1 + \tan^2 \theta = \sec^2 \theta$.
  \item Rearrange the identity to isolate $\sec^2 \theta - 1$.
\end{enumerate}
\textbf{Derivation:}
\begin{center}
\fitmath{\begin{aligned}
1 + \tan^2 \theta = \sec^2 \theta \\
\implies \sec^2 \theta - 1 = \tan^2 \theta
\end{aligned}}
\end{center}
\hfill \textit{Introduction to Trigonometry — Trigonometric Identities}
\vspace{0.3cm}

\question \textbf{36.} Assertion (A): If $\sin \theta = \frac{1}{2}$, then $\theta = 30^\circ$. Reason (R): The value of $\sin \theta$ increases as $\theta$ increases from $0^\circ$ to $90^\circ$.

\textbf{Answer:}
(B)

\textbf{Solution:}
\begin{enumerate}
  \item Check the validity of Assertion: For $\sin \theta = 1/2$, $\theta = 30^\circ$ is true.
  \item Check the validity of Reason: The sine function is strictly increasing in the first quadrant, which is true.
  \item Determine if R explains A: The fact that $\sin 30^\circ = 1/2$ is a specific value, its increasing nature is a general property, so R does not explain why $\sin 30^\circ$ is specifically $1/2$.
\end{enumerate}
\textbf{Derivation:}
\begin{center}
\fitmath{\begin{aligned}
\text{Assertion: } \sin 30^\circ = 0.5 \\
\implies \text{True. Reason: } \frac{d}{d\theta}(\sin \theta) = \cos \theta > 0 \text{ for } \theta \in (0, 90^\circ) \\
\implies \text{True.}
\end{aligned}}
\end{center}
\hfill \textit{Introduction to Trigonometry — Trigonometric Ratios of Specific Angles}
\vspace{0.3cm}

\question \textbf{37.} Assertion (A): The value of $\cos \theta$ can never be $1.5$. Reason (R): The range of $\cos \theta$ is $[-1, 1]$.

\textbf{Answer:}
(A)

\textbf{Solution:}
\begin{enumerate}
  \item Check the validity of Assertion: Since the range of cosine is $[-1, 1]$, $1.5$ is outside this range, so it is true.
  \item Check the validity of Reason: By definition of trigonometric ratios in a right triangle, $\cos \theta = \text{base/hypotenuse} \le 1$.
  \item Determine if R explains A: The range definition is the direct reason why $\cos \theta$ cannot exceed 1.
\end{enumerate}
\textbf{Derivation:}
\begin{center}
\fitmath{-1 \le \cos \theta \le 1 \text{ for all } \theta \in \mathbb{R}}
\end{center}
\hfill \textit{Introduction to Trigonometry — Basic Trigonometric Ratios}
\vspace{0.3cm}

\question \textbf{38.} Assertion (A): $\sin^2 \theta + \cos^2 \theta = 1$ for all $\theta$. Reason (R): $\tan^2 \theta - \sec^2 \theta = 1$.

\textbf{Answer:}
(C)

\textbf{Solution:}
\begin{enumerate}
  \item Check the validity of Assertion: This is the fundamental Pythagorean identity, which is true.
  \item Check the validity of Reason: The identity is $\sec^2 \theta - \tan^2 \theta = 1$. Therefore, $\tan^2 \theta - \sec^2 \theta = -1$. The reason is false.
\end{enumerate}
\textbf{Derivation:}
\begin{center}
\fitmath{\begin{aligned}
\sin^2 \theta + \cos^2 \theta = 1 \text{ (Identity)}; \sec^2 \theta - \tan^2 \theta = 1 \\
\implies \tan^2 \theta - \sec^2 \theta = -1
\end{aligned}}
\end{center}
\hfill \textit{Introduction to Trigonometry — Trigonometric Identities}
\vspace{0.3cm}

\question \textbf{39.} Assertion (A): If $\tan A = \cot A$, then $A = 45^\circ$. Reason (R): $\tan A = \cot(90^\circ - A)$.

\textbf{Answer:}
(A)

\textbf{Solution:}
\begin{enumerate}
  \item Check the validity of Assertion: If $\tan A = \cot A$, then $\tan A = 1/\tan A$, so $\tan^2 A = 1$, $\tan A = 1$, $A = 45^\circ$. True.
  \item Check the validity of Reason: This is the standard co-function identity.
  \item Determine if R explains A: Yes, $\tan A = \cot A \implies \tan A = \tan(90^\circ - A) \implies A = 90^\circ - A \implies 2A = 90^\circ \implies A = 45^\circ$.
\end{enumerate}
\textbf{Derivation:}
\begin{center}
\fitmath{\begin{aligned}
\tan A = \cot A \\
\iff \tan A = \tan(90^\circ - A) \\
\iff A = 90^\circ - A \\
\iff A = 45^\circ
\end{aligned}}
\end{center}
\hfill \textit{Introduction to Trigonometry — Trigonometric Ratios of Complementary Angles}
\vspace{0.3cm}

\question \textbf{40.} Assertion (A): The value of $\sin 0^\circ$ is $0$. Reason (R): The value of $\cos 0^\circ$ is $0$.

\textbf{Answer:}
(C)

\textbf{Solution:}
\begin{enumerate}
  \item Check the validity of Assertion: $\sin 0^\circ = 0$. True.
  \item Check the validity of Reason: $\cos 0^\circ = 1$. False.
\end{enumerate}
\textbf{Derivation:}
\begin{center}
\fitmath{\sin 0^\circ = 0, \cos 0^\circ = 1}
\end{center}
\hfill \textit{Introduction to Trigonometry — Trigonometric Ratios of Specific Angles}
\vspace{0.3cm}

\end{questions}

\end{document}
